\begin{helper}[The marks sit on their own time level]
Let $F$ be a finite-dimensional real inner-product space, let
$C_0,\kappa,H\in\R$, let $v$ be a real function on space-time, let
$T\in\mathbb N$, $R\in\R$, $N\in\mathbb N$ and $\epsilon\in\R$, and let
$(\Omega,\mathcal M,\mu)$ be a measure space with a filtration
$(\mathcal G_n)_{n\in\mathbb N}$ carrying functions $\xi_n$, $z_n$, $z^+_n$,
$z^-_n$, $Q_0(n)$ and $\tau$ which form a stopped contact walk for
$(\kappa,H,T,R,v,N,\epsilon)$ in the sense of \noderef{IsContactWalk}, together
with functions $\sigma_j$, $p_j$, $a_j$ which are mark data for
$(C_0,v,N,\mathcal G,z,\tau)$ in the sense of \noderef{IsMarkData}. Write
$\delta_i=2^{-i}$ and $z_n=(x_n,t_n)$. Then for every $j$ with $1\le j\le N$:
\begin{enumerate}
\item[(i)] $t_{\sigma_j}\ge\delta_j$ everywhere on $\Omega$;
\item[(ii)] $t_{\sigma_j}=\delta_j$ at every point where $t_\tau=\delta_N$.
\end{enumerate}
\end{helper}
\begin{proof}
Clause numbers refer to \noderef{IsContactWalk} unless \noderef{IsMarkData} is
named. Fix $j$ with $1\le j\le N$.

\emph{(G1) The grid.} By clause 2 there is an integer $k\ge0$ with
$\epsilon=2^{-k}$, and $0<\epsilon\le\delta_N/4=2^{-N-2}$; since $i\mapsto2^{-i}$
is decreasing, $k\ge N+2$. Consequently, for every $i\le k$ the number
$\delta_i=2^{-i}$ is an integer multiple of $\epsilon$, namely
$\delta_i=2^{\,k-i}\epsilon$, because $2^{\,k-i}2^{i}=2^k$ and
$2^{-i}2^{i}=1$. In particular $1=\delta_0=2^k\epsilon$ and
$\delta_j=2^{\,k-j}\epsilon$ are integer multiples of $\epsilon$, using
$j\le N\le k$.

\emph{(G2) The time coordinate moves by $0$ or $\pm\epsilon$.} Let $n$ and
$\omega$ be given. If $n\ge\tau(\omega)$ then clause 7 gives
$z_{n+1}=z_n$, so $t_{n+1}=t_n$. If $n<\tau(\omega)$ then clause 6 makes
$(z^+_n,z^-_n)$ an optimizing step pair of size $\epsilon$ at $z_n$, so by
\noderef{IsStepPair} the time coordinates of $z^+_n$ and $z^-_n$ are
$t_n+\epsilon$ and $t_n-\epsilon$; clause 7 says that $z_{n+1}$ is one of the
two, so $t_{n+1}\in\{t_n+\epsilon,\;t_n-\epsilon\}$. In all cases
$t_{n+1}\in\{t_n,\;t_n+\epsilon,\;t_n-\epsilon\}$.

\emph{(G3) Every time coordinate is an integer multiple of $\epsilon$.} By
induction on $n$. For $n=0$, clause 5 gives $z_0=(0,1)$, so $t_0=1$, which is an
integer multiple of $\epsilon$ by (G1). If $t_n=m\epsilon$ with $m\in\mathbb Z$
then by (G2) $t_{n+1}$ is one of $m\epsilon$, $(m+1)\epsilon$, $(m-1)\epsilon$.

\emph{(G4) Two multiples of $\epsilon$ in a half-open interval of length
$\epsilon$ coincide.} Let $x=m_1\epsilon$ and $c=m_2\epsilon$ with
$m_1,m_2\in\mathbb Z$ and $c-\epsilon<x\le c$. Dividing by $\epsilon>0$ gives
$m_2-1<m_1\le m_2$ in $\R$, hence in $\mathbb Z$, so $m_1=m_2$ and $x=c$.

\emph{(G5) A discrete intermediate value property.} Let $\omega\in\Omega$, let
$c$ be an integer multiple of $\epsilon$, let $b\in\mathbb N$ and suppose
$t_b(\omega)\le c\le t_0(\omega)$. Then $t_n(\omega)=c$ for some $n\le b$. We
induct on $b$. For $b=0$ we have $t_0(\omega)\le c\le t_0(\omega)$, so
$t_0(\omega)=c$. For the step, suppose the claim holds for $b$ and let
$t_{b+1}(\omega)\le c\le t_0(\omega)$. If $t_b(\omega)\le c$ the inductive
hypothesis gives $n\le b\le b+1$ with $t_n(\omega)=c$. Otherwise
$t_{b+1}(\omega)\le c<t_b(\omega)$; of the three possibilities of (G2) only
$t_{b+1}(\omega)=t_b(\omega)-\epsilon$ is compatible with
$t_{b+1}(\omega)\le c<t_b(\omega)$, so
$t_{b+1}(\omega)=t_b(\omega)-\epsilon>c-\epsilon$. Thus $t_{b+1}(\omega)$ and
$c$ are integer multiples of $\epsilon$ --- by (G3) and the hypothesis on $c$ ---
lying in $(c-\epsilon,c]$, so $t_{b+1}(\omega)=c$ by (G4), and $n=b+1$ serves.

\emph{(i).} Note $\delta_j\le\delta_0=1=t_0(\omega)$ for every $\omega$, by
clause 5. Suppose $t_{\sigma_j}(\omega)<\delta_j$ for some $\omega$. Since
$\delta_j$ is an integer multiple of $\epsilon$ by (G1) and
$t_{\sigma_j(\omega)}(\omega)\le\delta_j\le t_0(\omega)$, (G5) applied with
$b=\sigma_j(\omega)$ and $c=\delta_j$ gives an $n\le\sigma_j(\omega)$ with
$t_n(\omega)=\delta_j$. As $t_{\sigma_j}(\omega)<\delta_j$ we have
$n\ne\sigma_j(\omega)$, so $n<\sigma_j(\omega)$; but then clause 3 of
\noderef{IsMarkData} gives $t_n(\omega)\ne\delta_j$, a contradiction. Hence
$t_{\sigma_j}\ge\delta_j$ everywhere.

\emph{(ii).} Let $\omega$ satisfy $t_\tau(\omega)=\delta_N$. Since $j\le N$ we
have $\delta_N\le\delta_j$, so
$t_{\tau(\omega)}(\omega)=\delta_N\le\delta_j\le1=t_0(\omega)$, and (G5) with
$b=\tau(\omega)$ and $c=\delta_j$ gives an $n\le\tau(\omega)$ with
$t_n(\omega)=\delta_j$. Let $n_0$ be the least index with
$t_{n_0}(\omega)=\delta_j$; then $n_0\le n\le\tau(\omega)$. If
$n_0<\sigma_j(\omega)$ then clause 3 of \noderef{IsMarkData} applied at $n_0$
would give $t_{n_0}(\omega)\ne\delta_j$; hence $\sigma_j(\omega)\le n_0$. By
clause 4 of \noderef{IsMarkData} either $t_{\sigma_j}(\omega)=\delta_j$, which
is the claim, or $\sigma_j(\omega)=\tau(\omega)$, in which case
$n_0\le\tau(\omega)=\sigma_j(\omega)\le n_0$ forces
$\sigma_j(\omega)=n_0$ and therefore $t_{\sigma_j}(\omega)=\delta_j$ again.
\end{proof}
